% References used in the revised narrative; original reference list is archived separately.
\setlength{\bibsep}{3pt plus 1pt minus 1pt}
\begin{thebibliography}{00}
\bibitem{berecibar2016} M. Berecibar, I. Gandiaga, I. Villarreal, N. Omar, J. Van Mierlo, P. Van den Bossche, Critical review of state of health estimation methods of Li-ion batteries for real applications, Renewable and Sustainable Energy Reviews 56 (2016) 572--587.
\bibitem{waag2014} W. Waag, C. Fleischer, D.U. Sauer, Critical review of the methods for monitoring of lithium-ion batteries in electric and hybrid vehicles, Journal of Power Sources 258 (2014) 321--339.
\bibitem{farmann2015} A. Farmann, W. Waag, A. Marongiu, D.U. Sauer, Critical review of on-board capacity estimation techniques for lithium-ion batteries in electric and hybrid electric vehicles, Journal of Power Sources 281 (2015) 114--130.
\bibitem{ecker2012} M. Ecker, J.B. Gerschler, J. Vogel, S. K{\"o}bitz, F. Hust, P. Dechent, D.U. Sauer, Development of a lifetime prediction model for lithium-ion batteries based on extended accelerated aging test data, Journal of Power Sources 215 (2012) 248--257.
\bibitem{waldmann2014} T. Waldmann, M. Wilka, M. Kasper, M. Fleischhammer, M. Wohlfahrt-Mehrens, Temperature dependent ageing mechanisms in lithium-ion batteries--A post-mortem study, Journal of Power Sources 262 (2014) 129--135.
\bibitem{dubarry2012} M. Dubarry, C. Truchot, B.Y. Liaw, Synthesize battery degradation modes via a diagnostic and prognostic model, Journal of Power Sources 219 (2012) 204--216.
\bibitem{rasmussen2006} C.E. Rasmussen, C.K.I. Williams, Gaussian Processes for Machine Learning, MIT Press, 2006.
\bibitem{cortes1995} C. Cortes, V. Vapnik, Support-vector networks, Machine Learning 20 (1995) 273--297.
\bibitem{breiman2001} L. Breiman, Random forests, Machine Learning 45 (2001) 5--32.
\bibitem{chen2016} T. Chen, C. Guestrin, XGBoost: A scalable tree boosting system, Proceedings of the 22nd ACM SIGKDD International Conference (2016) 785--794.
\bibitem{severson2019} K.A. Severson et al., Data-driven prediction of battery cycle life before capacity degradation, Nature Energy 4 (2019) 383--391.
\bibitem{pan2010} S.J. Pan, Q. Yang, A survey on transfer learning, IEEE Transactions on Knowledge and Data Engineering 22 (2010) 1345--1359.
\bibitem{long2015} M. Long, Y. Cao, J. Wang, M. Jordan, Learning transferable features with deep adaptation networks, Proceedings of the 32nd International Conference on Machine Learning (2015) 97--105.
\bibitem{sun2016} B. Sun, K. Saenko, Deep CORAL: correlation alignment for deep domain adaptation, European Conference on Computer Vision Workshops (2016) 443--450.
\bibitem{ganin2016} Y. Ganin, E. Ustinova, H. Ajakan, P. Germain, H. Larochelle, F. Laviolette, M. Marchand, V. Lempitsky, Domain-adversarial training of neural networks, Journal of Machine Learning Research 17 (2016) 1--35.
\bibitem{mackay1998} D.J.C. MacKay, Introduction to Gaussian processes, in: C.M. Bishop (Ed.), Neural Networks and Machine Learning, Springer, 1998, pp. 133--165.
\bibitem{calce_gp2013} D. Liu, J. Pang, J. Zhou, Y. Peng, M. Pecht, Prognostics for state of health estimation of lithium-ion batteries based on combination Gaussian process functional regression, Microelectronics Reliability 53 (6) (2013) 832--839.
\bibitem{wang2024} F. Wang, Z. Zhai, Z. Zhao, Y. Di, X. Chen, Physics-informed neural network for lithium-ion battery degradation stable modeling and prognosis, Nature Communications 15 (2024) 4332.
\bibitem{lu2023} J. Lu, R. Xiong, J. Tian, C. Wang, F. Sun, Deep learning to estimate lithium-ion battery state of health without additional degradation experiments, Nature Communications 14 (2023) 2760.
\bibitem{zhu2022} J. Zhu et al., Data-driven capacity estimation of commercial lithium-ion batteries from voltage relaxation, Nature Communications 13 (2022) 2261.
\bibitem{batteryml2024} H. Zhang, X. Gui, S. Zheng, Z. Lu, Y. Li, J. Bian, BatteryML: An Open-source Platform for Machine Learning on Battery Degradation, International Conference on Learning Representations, 2024. Software and data processing: \url{https://github.com/microsoft/BatteryML}.
\bibitem{batterylife2025} R. Tan, W. Hong, J. Tang, X. Lu, R. Ma, X. Zheng, J. Li, J. Huang, T.-Y. Zhang, BatteryLife: A Comprehensive Dataset and Benchmark for Battery Life Prediction, Proceedings of the 31st ACM SIGKDD Conference on Knowledge Discovery and Data Mining V.2 (2025) 5789--5800. Processed-data release v12: \url{https://github.com/Ruifeng-Tan/BatteryLife}.
\bibitem{nie2023} Y. Nie, N.H. Nguyen, P. Sinthong, J. Kalagnanam, A Time Series is Worth 64 Words: Long-term Forecasting with Transformers, International Conference on Learning Representations, 2023. \url{https://arxiv.org/abs/2211.14730}.
\bibitem{bai2018} S. Bai, J.Z. Kolter, V. Koltun, An Empirical Evaluation of Generic Convolutional and Recurrent Networks for Sequence Modeling, arXiv:1803.01271 (2018). \url{https://arxiv.org/abs/1803.01271}.
\bibitem{finn2017} C. Finn, P. Abbeel, S. Levine, Model-Agnostic Meta-Learning for Fast Adaptation of Deep Networks, Proceedings of Machine Learning Research 70 (2017) 1126--1135. \url{https://proceedings.mlr.press/v70/finn17a.html}.
\bibitem{zhao2025} X. Zhao, Z. Wang, T. Han, W. Yang, F. Gu, A.D. Ball, A Meta-Learning Method for Few-Shot Multidomain State-of-Health Estimation of Lithium-Ion Batteries, IEEE Transactions on Transportation Electrification 11 (1) (2025) 4830--4840.
\bibitem{ulpur2020} D. Juarez-Robles, A.A. Vyas, C. Fear, J.A. Jeevarajan, P.P. Mukherjee, Overdischarge and Aging Analytics of Li-Ion Cells, Journal of The Electrochemical Society 167 (9) (2020) 090558. UL-PUR data resource: \url{https://ul.org/research/electrochemical-safety/open-science-data/aging-and-safety-pouch-cells}.
\bibitem{mich2021} A. Weng, P. Mohtat, P.M. Attia, V. Sulzer, S. Lee, G. Less, A. Stefanopoulou, Predicting the impact of formation protocols on battery lifetime immediately after manufacturing, Joule 5 (11) (2021) 2971--2992. MICH data repository: University of Michigan Deep Blue Data.
\bibitem{stroe2018} D.-I. Stroe, V. Knap, E. Schaltz, State-of-Health Estimation of Lithium-Ion Batteries based on Partial Charging Voltage Profiles, ECS Transactions 85 (13) (2018) 379--386.
\bibitem{meng2019} J. Meng, L. Cai, D.-I. Stroe, G. Luo, X. Sui, R. Teodorescu, Lithium-ion battery state-of-health estimation in electric vehicle using optimized partial charging voltage profiles, Energy 185 (2019) 1054--1062.
\bibitem{tang2024} J. Tang, Y. Li, S. Wang, B. Xiong, X. Li, J. Pan, Q. Chen, P. Wang, Data-driven state of health estimation method of lithium-ion batteries for partial charging curves, IEEE Transactions on Energy Conversion 39 (4) (2024) 2230--2243.
\bibitem{xiong2023} X. Xiong, Y. Wang, K. Li, Z. Chen, State of health estimation for lithium-ion batteries using Gaussian process regression-based data reconstruction method during random charging process, Journal of Energy Storage 72 (2023) 108390.
\bibitem{lyu2020} Z. Lyu, R. Gao, Li-ion battery state of health estimation through Gaussian process regression with Thevenin model, International Journal of Energy Research 44 (13) (2020).
\bibitem{yu2018} J. Yu, State of health prediction of lithium-ion batteries: Multiscale logic regression and Gaussian process regression ensemble, Reliability Engineering \& System Safety 174 (2018) 82--95.
\bibitem{jin2023} H. Jin, N. Cui, L. Cai, J. Meng, J. Li, J. Peng, X. Zhao, State-of-health estimation for lithium-ion batteries with hierarchical feature construction and auto-configurable Gaussian process regression, Energy 262 (2023) 125503.
\end{thebibliography}
